\documentclass[11pt,a4paper]{article}
\usepackage[margin=1.8cm]{geometry}
\usepackage{amsmath,amssymb,tabularx,array,needspace,fontspec,enumitem,cancel}
\newcommand{\tb}{\cancel{||||}\,}
\setmainfont{FreeSerif}
\newfontfamily\urdufont[Script=Arabic]{FreeSerif}
\TeXXeTstate=1
\newcommand{\urdu}[1]{{\urdufont\beginR #1\endR}}
\setlength{\parindent}{0pt}
\setlength{\parskip}{0pt}
\renewcommand{\arraystretch}{1.15}
\pagestyle{plain}

% \mcq{number}{question}{A}{B}{C}{D}
\newcommand{\mcq}[6]{%
  \Needspace{4\baselineskip}%
  \par\medskip\noindent\textbf{#1.}\ #2\par\nopagebreak
  \noindent\begin{tabularx}{\linewidth}{@{\hspace{1.5em}}*{4}{>{\raggedright\arraybackslash}X}@{}}
    (A)~#3 & (B)~#4 & (C)~#5 & (D)~#6
  \end{tabularx}\par}

\begin{document}

\begin{center}
  {\Huge\bfseries The Hope Science Academy}\\[3pt]
  {\small Excellence is our standard not goal}\\[2pt]
  {\small 247-E-1, (Opposite LDA School) Johar Town, Lahore. PH: 0300-8194789, 0322-4157001}
\end{center}

\medskip
\noindent
\begin{tabularx}{\linewidth}{|X|X|X|X|}
\hline
\textbf{Student Name:} & \textbf{Roll Num:} & \textbf{Class:} 11th & \textbf{Paper Code:} \\ \hline
\textbf{Subject Name:} Business Maths & \textbf{Time Allowed:} & \textbf{Total Marks:} 490 & \textbf{Exam Date:} 02-Nov-2026 \\ \hline
\multicolumn{4}{|l|}{\textbf{Exam Syllabus:} FULL BOOK} \\ \hline
\end{tabularx}

\bigskip
\noindent\textbf{Q1. Choose the correct answer.}\hfill\textbf{1$\times$490=490}

\mcq{1}{The ratio between two homogenous quantities or number is called:}{Proportion}{Annuity}{Ratio}{Percentage}
\mcq{2}{Comparison between two quantities of same unites in known as:}{Percentage}{Ratio}{Proportion}{Interest}
\mcq{3}{The ratio 168:147 in lowest form is:}{24;21}{56;49}{8;7}{7;8}
\mcq{4}{18 to 30 is same as:}{3:4}{3:5}{6:7}{6:9}
\mcq{5}{Reduced form of 36;60:}{12:5}{3:5}{1:2}{3:10}
\mcq{6}{The ration between 80 and 640:}{1:4}{2:4}{1:8}{4:16}
\mcq{7}{Ratio between 10 minutes and 30 minutes is:}{2:3}{1:3}{2:4}{1:5}
\mcq{8}{15 seconds 3 minutes =:}{1:12}{5:1}{1:5}{12:1}
\mcq{9}{The ration betwen 3.5 kg and 7 kg is:}{15:7}{7:15}{7:5}{5:15}
\mcq{10}{The ratio between 7.5 kg and 3.5 kg:}{15:7}{7:15}{7:15}{5:15}
\mcq{11}{1 kg ; 700 grams is:}{10:7}{7:10}{1:700}{700:1}
\mcq{12}{5 to 25 is same as:}{$5:\dfrac{1}{25}$}{$\dfrac{1}{5}:25$}{5:25}{None}
\mcq{13}{3.5 is a ratio of:}{3 and 25}{3.5 and 100}{325 and 100}{13 and 40}
\mcq{14}{The ratio $\dfrac{4}{9}:\dfrac{1}{3}$ in lowest term is:}{4:3}{3:4}{9:3}{3:1}
\mcq{15}{The ratio between 2.5 kg and 4.5 kg is:}{5:9}{2:5}{9:5}{5:2}
\mcq{16}{The ratio between 75 and 225 in its lowest term is:}{1:3}{2:4}{1:8}{4:16}
\mcq{17}{$\dfrac{5}{6}$ to $\dfrac{10}{9}$ is same as:}{5:9}{5:10}{3:4}{3:5}
\mcq{18}{The ratio between 1.9 kg and 9.5 kg is:}{1:5}{5:1}{1:4}{4:1}
\mcq{19}{The decimal fraction of 16.35\%=......:}{0.1635}{1.635}{163.5}{1635.0}
\mcq{20}{Decimal form of 5.3\% is:}{53}{0.53}{0.0053}{0.053}
\mcq{21}{Reduced form of ratio 27:54 is:}{2:1}{1:2}{1:3}{3:1}
\mcq{22}{The simplest ratio of 18 to 30 is:}{3:4}{3:5}{6:9}{6:7}
\mcq{23}{Reduce 40;240 to its simplest form:}{1:6}{1:3}{1:4}{None of these}
\mcq{24}{Reduce form of ratio 16 to 64 is:}{1:4}{4:1}{4:3}{3:4}
\mcq{25}{If profit distribution ratio of two partners is 2 to 3 then second partner's share in profit is:}{$\dfrac{2}{3}$}{$\dfrac{3}{2}$}{$\dfrac{3}{5}$}{$\dfrac{2}{5}$}
\mcq{26}{$\dfrac{2}{9}:\dfrac{4}{36}:\dfrac{16}{18}=$ :}{4:3:2}{2:3:4}{2:1:8}{1:3:4}
\mcq{27}{The 1:3 is the reduced form of:}{14:34}{13:33}{12:32}{11:33}
\mcq{28}{The ratio between 850 grams and 1.70 kg is:}{5:1}{1:2}{1:5}{2:1}
\mcq{29}{Reduced form of 27:54 is:}{2:1}{3:4}{3:1}{1:2}
\mcq{30}{The ratio between 80 and 640 is:}{1:4}{2:4}{1:8}{4:16}
\mcq{31}{3.25 is a ratio of:}{3 and 25}{32.5 and 10}{325 and 100}{13 and 4}
\mcq{32}{The ratio between 7.5 kg and 3.5 kg is:}{7:15}{7 to 15}{7:25}{None}
\mcq{33}{If A and B are in the business and A's investment in business is $\dfrac{a}{a+b}$, then B's investment would be:}{$\dfrac{a+b}{a}$}{$\dfrac{b}{a+b}$}{$\dfrac{a+b}{b}$}{$1-\dfrac{b}{a+b}$}
\mcq{34}{A statement expressing equality of two ratio is called:}{Ratio}{Proportion}{Percentage}{Annuity}
\mcq{35}{The missing quantity in the following x;5;;12;30 is:}{2}{12}{3}{1}
\mcq{36}{If 3;12;;180;x then x is equal to:}{420}{520}{620}{720}
\mcq{37}{The missing quantity is the statement x;5;;15;25 is:}{3}{5}{15}{25}
\mcq{38}{f 2;7;;x;49 then the value of x is:}{9}{14}{52}{26}
\mcq{39}{If 7;x=7;8 then value of x is:}{7}{1}{56}{8}
\mcq{40}{5;10;;12;24 the proportion are:}{Inverse}{Direct}{Continued}{Compound}
\mcq{41}{a;3;;b;9 is same as:}{$\dfrac{a}{3}=\dfrac{b}{3}$}{$\dfrac{a}{1}=\dfrac{b}{3}$}{$\dfrac{a}{a}=\dfrac{b}{27}$}{$\dfrac{a}{3}=\dfrac{b}{27}$}
\mcq{42}{$\dfrac{40}{6}=\dfrac{80}{x}$ is equivalent to:}{7:4::8:30}{4:7::8:30}{4:3::8:7}{7:8::4:30}
\mcq{43}{If x;4;;3;2 then the value of x is:}{2}{4}{8}{6}
\mcq{44}{If a;b=c;d= them:}{ab=cd}{ad=bc}{ac=bd}{None of these}
\mcq{45}{If 39;x;;18;6, then x is equal to:}{15}{13}{17}{19}
\mcq{46}{In 10;x ;; 8;4 is equal to:}{5}{6}{4}{2}
\mcq{47}{The missing value in $\dfrac{40}{6}=\dfrac{80}{x}$ is:}{12}{14}{3}{None of these}
\mcq{48}{The missing value in $\dfrac{20}{7}=\dfrac{80}{x}$ is:}{2}{24}{26}{28}
\mcq{49}{If $\dfrac{3}{15}=\dfrac{7}{x}$ then x is equal to:}{25}{35}{30}{40}
\mcq{50}{If price of three shirts is Rs. 750, then the price of two shirts will be:}{500}{600}{700}{800}
\mcq{51}{The price of 12 bananas is Rs.50. What is the price of 24 bananas:}{100}{60}{700}{80}
\mcq{52}{$\dfrac{a}{b}=\dfrac{b}{c}=\dfrac{c}{d}=\dfrac{d}{e}$ is an example of:}{Direct proportion}{Inverse proportion}{Compound proportion}{Continued proportion}
\mcq{53}{If A;B= 2;3 and B;C=3;5, then A;B;C would be:}{2:3:5}{2:5:3}{3:2:5}{3:5:2}
\mcq{54}{If A;B= 3;5 and B;C=5;6 hence A;B;C would be:}{5:3:6}{3:5:6}{3:6:5}{6:5:3}
\mcq{55}{The middle term of every proportion are called:}{Central terms}{Extremes}{Means}{Upper terms}
\mcq{56}{The upper and lower terms as regard to position of a proportion are called:}{Means}{Extremes}{Highest \& lowest terms}{Focal terms}
\mcq{57}{If 50 persons made 500 tables in 5 days. The rate of production per worker per day is:}{1 table}{2 tables}{3 tables}{4 tables}
\mcq{58}{$\dfrac{35}{125}=\dfrac{7}{x}$, then x is:}{25}{35}{30}{40}
\mcq{59}{0.047 is equivalent to:}{0.047\%}{0.47\%}{4.7\%}{47\%}
\mcq{60}{The change of 0.2\% to fractional is:}{$\dfrac{1}{100}$}{$\dfrac{1}{200}$}{$\dfrac{1}{400}$}{$\dfrac{1}{500}$}
\mcq{61}{25\% means:}{$\dfrac{3}{4}$}{$\dfrac{1}{4}$}{$\dfrac{1}{2}$}{$\dfrac{1}{6}$}
\mcq{62}{0.00425 is percentage form is:}{4.25\%}{0.425\%}{42.5\%}{425\%}
\mcq{63}{The decimal fraction of 16.35\% =......:}{0.1635}{1.635}{163.5}{1635.0}
\mcq{64}{\% of 400 is:}{20}{40}{60}{80}
\mcq{65}{20\% of 500 is:}{50}{100}{150}{200}
\mcq{66}{25\% os 500 is:}{100}{200}{150}{125}
\mcq{67}{30\% of 300 is:}{80}{90}{70}{60}
\mcq{68}{What number reduced by 25\% of itself becomes 165:}{200}{220}{240}{280}
\mcq{69}{30\% of 1800 is:}{450}{540}{240}{340}
\mcq{70}{5 out of 50 students failed, what is pass percentage:}{80\%}{75\%}{90\%}{95\%}
\mcq{71}{What percent Rs. 30 is of 300:}{20\%}{10\%}{5\%}{2\%}
\mcq{72}{160 is 20\% of what numbr:}{80}{800}{8000}{80000}
\mcq{73}{10 out 50 candiate fail in exam. What is fail \%:}{10}{20}{30}{40}
\mcq{74}{What percent of Rs. 20 is of 200:}{20\%}{10\%}{15\%}{25\%}
\mcq{75}{30\% of 12000 is:}{30}{300}{3000}{3600}
\mcq{76}{$\dfrac{1}{2}$ is what percentage of $\dfrac{1}{4}$:}{25\%}{50\%}{100\%}{200\%}
\mcq{77}{20\% of 1250 is:}{25}{250}{2500}{625}
\mcq{78}{195 is what percentage of 78:}{40\%}{200\%}{250\%}{100\%}
\mcq{79}{45\% of 800 is:}{350}{315}{355}{360}
\mcq{80}{Asghar spends 80\% of income and save Rs. 1200. What is total income of Asghar:}{RS. 5700}{Rs.5800}{Rs.5900}{Rs.6000}
\mcq{81}{Rs 250 is what percent of 1000:}{1.5}{2.5}{3.5}{None of these}
\mcq{82}{What percent Rs. 45 is of 250:}{16\%}{17\%}{18\%}{19\%}
\mcq{83}{150 is 25\% of number:}{600}{60000}{60}{6000}
\mcq{84}{27 is what percentage of 300:}{3\%}{6\%}{9\%}{12\%}
\mcq{85}{30\% of 900 is:}{150}{250}{270}{370}
\mcq{86}{Rs 250 is $2\dfrac{1}{2}$\% of Rs.:}{10000}{1000}{100}{10}
\mcq{87}{20\% of 1600 is:}{3200}{320}{80}{32}
\mcq{88}{What percentage of 10000 is 525:}{0.525\%}{5.25\%}{52.5\%}{15\%}
\mcq{89}{25 is what percentage of 500:}{15\%}{10\%}{5\%}{20\%}
\mcq{90}{10\% of 300 is:}{10}{30}{100}{150}
\mcq{91}{250 is 20\% of:}{1250}{1230}{1150}{1050}
\mcq{92}{What percent Rs 88 is of Rs. 400:}{22\%}{21\%}{23\%}{25\%}
\mcq{93}{...... of 560 is 28:}{20\%}{15\%}{10\%}{5\%}
\mcq{94}{20\% is what percent of Rs.80:}{15\%}{20\%}{25\%}{30\%}
\mcq{95}{Decimal form of 5.3\% is:}{53}{0.53}{0.0053}{0.053}
\mcq{96}{The equality of two ratio is called:}{Proportion}{Direct proportion}{Indirect proportion}{Percentage}
\mcq{97}{The price before deducting the discount is called:}{Gross price}{List price}{Cost price}{Net price}
\mcq{98}{Number of types of discount are:}{Four}{Three}{Two}{One}
\mcq{99}{The amount after deduction trade discount and cash discount is called:}{Cost price}{Selling price}{Net price}{None}
\mcq{100}{The discount is:}{Reduction}{Increase}{May increase of decrease}{None of these}
\mcq{101}{Buy 3 get 5 offer contains:}{Discount @ 66.67\%}{Discount @ 33.33\%}{Discount @ 40\%}{Discount at unknown rate}
\mcq{102}{Formula for finding discounted price with three successive discounts is:}{SP=DP(1+DR$_{1}$)(1-DR$_{2}$)(1-DR$_{3}$)}{DP=SP(1-DR$_{1}$)(1-DR$_{2}$)(1+DR$_{3}$)}{DP=SP(1-DR$_{1}$)(1-DR$_{2}$)(1-DR$_{3}$)}{SP=DP(1-DR$_{1}$)(1-DR$_{2}$)(1-DR$_{3}$)}
\mcq{103}{The amount paid to an agent as remuneration for his services is called:}{Salary}{Perpetuity}{Commission}{Profit}
\mcq{104}{Total amount of Rs. 700,000 then commission @ 2\% is:}{14000}{12000}{10000}{21000}
\mcq{105}{Commission is:}{Remuneration}{Salary}{Wages}{None of these}
\mcq{106}{The percentage involved in brokerage is:}{Markup}{Profit rate}{Rate of commission}{None of these}
\mcq{107}{$2\dfrac{1}{2}$\% commission on a deal of Rs 150000 is:}{3750}{375.0}{375}{37500}
\mcq{108}{On a deal of Rs. 200000 the commission was Rs.4500. The commission charged at the rate of:}{2\%}{$2\dfrac{1}{4}$\%}{$2\dfrac{1}{2}$\%}{$2\dfrac{3}{4}$\%}
\mcq{109}{An agent collects commission @ $3\dfrac{1}{2}$\% and receives Rs 1400. The total amount of deal is:}{Rs. four thousand}{Rs forty thousand}{Rs. four lakh}{Rs. four million}
\mcq{110}{$\dfrac{Profit}{C.P}\times100\%=\ldots\ldots$:}{Mark up}{Mark down}{Loss}{Sale price}
\mcq{111}{What is cost price when sale price 4400 and profit is 10\%:}{4040}{4500}{4000}{4200}
\mcq{112}{Loss is opposite to:}{Markup}{Profit}{Mark-down}{Quotient}
\mcq{113}{The difference of selling price and cost as to be positive as:}{Loss}{Mark-up}{Profit}{Margin}
\mcq{114}{According to text (S-C)-0 is:}{Loss}{Profit}{Mark-up}{Mark-down}
\mcq{115}{Interest is:}{Wages}{Salary}{Rent}{Comission}
\mcq{116}{In Islam interest is:}{Compeletely forbidden}{Partially forbidden}{Allowed}{Sin}
\mcq{117}{Interest is classified in classes:}{Five}{Four}{Three}{Two}
\mcq{118}{The method of finding interest are:}{One}{Two}{Three}{Four}
\mcq{119}{Simple interest ``T'' is:}{P}{Prt/}{SRi}{P}
\mcq{120}{Pricipal amount ramian constant for calculating of interest for every period in:}{Simple interest}{Compound interest}{Rent}{Investment}
\mcq{121}{The formula for finding rate for simple interest is:}{$\dfrac{S.I}{P\times I}$}{$\dfrac{S.I}{P\times N}$}{$\dfrac{P\times I}{S.I}$}{$\dfrac{P\times N}{S.I}$}
\mcq{122}{Principal =Rs. 5000, interest rate = 10\% period =1/2 years, interest=:}{1000}{500}{250}{100}
\mcq{123}{Simple interest on Rs. 500 @ 30\% in 2 years is:}{150}{300}{600}{900}
\mcq{124}{The simple interest on Rs. 700 borrowed for one year at the rate of one percent per annum is:}{7}{70}{700}{80}
\mcq{125}{The simple interest on Rs. 700 borrowed for one year @ 8\% p.a is:}{250}{275}{276.28}{280.40}
\mcq{126}{An amount will by doubled @ 20\% with simple interest in:}{5 years}{20 years}{40 years}{10 years}
\mcq{127}{Simple interest on Rs. 500 on 12\% annually in 2 years is:}{60}{180}{360}{120}
\mcq{128}{If total amount is Rs. 1444 and principal 1300 then simple interest is:}{300}{344}{44}{144}
\mcq{129}{Simple interest on Rs. 400@ 9\% annually in 2 years is:}{36}{360}{72}{720}
\mcq{130}{Capital remain same thorghout total period in:}{Annuity}{Perpetuity}{Compound interest}{Simple interest}
\mcq{131}{Amount is equal to:}{Principal - interest}{Principal + interest}{Interest - principal}{Principal x interest}
\mcq{132}{The simple interest on Rs. 500, invested for 6 months at 8\% per annum is:}{30}{20}{25}{40}
\mcq{133}{If total amount is paid Rs. 1344 and simple interest is 344 then principal is:}{1300}{1100}{1000}{344}
\mcq{134}{If simple interest on 15000 for 3 years is Rs. 900 the rate of interest is:}{4\%}{3\%}{5\%}{2\%}
\mcq{135}{Rs. 500 on the simple interest double itself in 5 years at the rate of:}{20\%}{30\%}{40\%}{50\%}
\mcq{136}{In simple interest capital remains:}{Constant}{Variable}{Change}{Re-invest}
\mcq{137}{Simple interest on Rs. 20000=@ 10\% in 4 years is:}{2000}{4000}{8000}{16000}
\mcq{138}{If an amount is doubled in 5 years, then the simple interest rate is:}{10\%}{15\%}{20\%}{25\%}
\mcq{139}{Simple interest on Rs. 15000 for 3 years at 2\% is:}{1000}{900}{1100}{1500}
\mcq{140}{Simple interest on Rs.2000 @ 5\% annually in 4 years is equal to:}{100}{200}{400}{600}
\mcq{141}{Simple interest on Rs. 4000 @ 5\% for two yars is:}{400}{200}{450}{300}
\mcq{142}{If amount 1000 is borrowed for five years and paid Rs. 1250 at the rate of:}{5\%}{10\%}{7\%}{9\%}
\mcq{143}{Simple interest on Rs. 1600 @ 5\% annually in 3 years is equal to:}{120}{160}{180}{240}
\mcq{144}{Simple interest on Rs. 10000 at the rate of 10\% in 10 years is:}{100}{1000}{10000}{100000}
\mcq{145}{Simple interst on a loan of Rs. 3000 for 2 years at 7\% is:}{220}{320}{420}{520}
\mcq{146}{Rs. 5000 itself in 5 years at the rate of:}{10\%}{20\%}{30\%}{None of these}
\mcq{147}{The formula of compound interest is:}{P\{1+i\}$^{n}$}{P[\{1+i\}$^{n}$+1]}{P[\{1+n\}$^{n}$-1]}{P[}
\mcq{148}{Compound amount formula is:}{P}{Prt}{P}{None of these}
\mcq{149}{The compound interest on Rs. 1000 for 2 year at 6\% is:}{120}{130}{124}{125}
\mcq{150}{An amount trippled itself @ 20\% in ...... years at compound interest is:}{2}{3}{4}{6}
\mcq{151}{Compound interest for same principal, same rate and sama period is always:}{$<$ simple interest}{$\leq$ Simple interest}{$>$ simple interest}{$\geq$ Simple interest}
\mcq{152}{A sequence of payment made at equal intervals of time is called:}{Perpetuity}{Annuity}{Compound amount}{None of these}
\mcq{153}{Regular, periodic and fixed sequence of payment is:}{Simple interest}{Annuity}{Compound interest}{Percentage}
\mcq{154}{The annuity in which terms of annuity is fixed:}{Ordinary annuity}{Annuity due}{Annuity certain}{Simple interest}
\mcq{155}{If every payment of an annuity is made at the beginning of th period, is called:}{Perpetuity}{Ordinary annuity}{Annuity due}{Compound amount}
\mcq{156}{Sum of annuity is called:}{Future value}{Present value}{Current value}{Past value}
\mcq{157}{Types of annuity are:}{Two}{Three}{Four}{Five}
\mcq{158}{Annuity is used in:}{Euqation}{Math of finance}{Linear function}{Quadratic function}
\mcq{159}{Monthly installment of a auto car on lease can be found by using the techmique of:}{Compound interest}{Annuity}{Simple interest method}{None of these}
\mcq{160}{$S=R\left[\dfrac{(1+i)^{n}-1}{i}\right]$ is formula for sum of amount of:}{Annuity due}{Ordinary annuity}{Simple interest}{Compound interest}
\mcq{161}{The time between successive payment for an annuity is called:}{Percentage}{Ratio}{Payment interval}{Total amount}
\mcq{162}{Reverse of sum of an annuity is called:}{Amount}{Present value}{Annuity}{Perpetuity}
\mcq{163}{Never ending annuity is:}{Ordinary annuity}{Perpetuity}{Annuity}{Annuity due}
\mcq{164}{The presnet value of an annuity:}{P}{S}{Sr}{A}
\mcq{165}{An ordinary annity which never stop is called:}{Annuity}{Perpetuity}{Present value}{Sum of annuity}
\mcq{166}{It is not possible to calculate:}{Sum of perpetuity}{Present value of ordinary annuity}{Compound interest}{Present value of annuity due}
\mcq{167}{Formula to calculate the present value of perpetuity is:}{$P=\dfrac{R}{i}$}{P=Ri}{$P=\dfrac{R}{S}$}{$P=\dfrac{i}{R}$}
\mcq{168}{Annuity classified into categories is:}{Four}{Three}{Two}{Five}
\mcq{169}{Sum of annuity is always:}{Future value}{Present value}{Current value}{None of these}
\mcq{170}{If the payment starts on a certain date and continuous for indefinite period, then it is called:}{Ordinary annuity}{Annuity due}{Perpetuity}{Contingent annuity}
\mcq{171}{In an equation, two expression are joined by the sign:}{$\leq$}{$\geq$}{$\neq$}{=}
\mcq{172}{An equation is a statement that two expression are:}{Equal}{Unequal}{May be equal}{None of these}
\mcq{173}{The standard form of linear equation:}{ax+b=0}{3c}{$4c^{2}+2c$}{None of these}
\mcq{174}{ax+by=c is an example of:}{Exponential equation}{Linear equation}{Identity equation}{Quadratic equation}
\mcq{175}{A first degree equaton is called:}{Non-linear}{Quadratic}{Constant}{Linear}
\mcq{176}{The solution set of equation ax+b=0 is:}{$\dfrac{a}{b}$}{$\dfrac{b}{a}$}{$\dfrac{-b}{a}$}{$\dfrac{-a}{b}$}
\mcq{177}{In equation 3x-6=0 the quantity x is called:}{Unknown}{Constant}{Coefficient}{None}
\mcq{178}{A linear equation has:}{Two roots}{Three roots}{One root}{None of these}
\mcq{179}{The equation ax+b=0 is:}{Linear}{Quadratic}{Constant}{None}
\mcq{180}{In the equation y=ax+b, a and b are:}{Variable}{Degree}{Constant number}{None}
\mcq{181}{Degree of linear equation is:}{One}{Two}{Three}{Four}
\mcq{182}{The degree of the equation 2x+5=0 is:}{2}{5}{1}{3}
\mcq{183}{The degree of an equation is highest power of:}{Co-efficient}{Constant}{Variable}{Exponents}
\mcq{184}{If three time a number is 150 the number is:}{25}{40}{50}{60}
\mcq{185}{If 2x-5=3 the value of x is:}{9}{6}{4}{16}
\mcq{186}{If 5 times a number is 190. than the number will be:}{40}{39}{38}{36}
\mcq{187}{If 5x+70=0 then x=......:}{12}{-13}{-14}{-15}
\mcq{188}{if x+9=15 then value of x is:}{6}{-6}{7}{24}
\mcq{189}{Which value of x is the root of the equation 11x-22=11:}{x=3}{x=-3}{x=4}{x=33}
\mcq{190}{If 8 times a number is 56 what is the number:}{5}{20}{6}{7}
\mcq{191}{Solution set of equation x-4=16 is:}{\{12\}}{\{20\}}{\{4\}}{\{-20\}}
\mcq{192}{Given that 3x+9=6x then:}{x=3}{x=4}{x=5}{x=6}
\mcq{193}{If six times a number is 180 then the number will be:}{10}{20}{30}{40}
\mcq{194}{If $\dfrac{1}{4}$ of an amount is 20, the total amount is:}{40}{80}{120}{16}
\mcq{195}{The solution set of x+9=2 then x is:}{\{9\}}{\{2\}}{$\left\{\dfrac{7}{2}\right\}$}{\{-7\}}
\mcq{196}{If 3x+2=2x+8 then value of x is:}{6}{5}{4}{3}
\mcq{197}{If 5 times a number is 115 then the number is:}{21}{27}{23}{19}
\mcq{198}{The equation y=7x+4 is called:}{Quadratic}{Linear}{Cubic}{Reciprocal}
\mcq{199}{The equation 3x-7=0 solution:}{$\dfrac{7}{3}$}{$\dfrac{3}{7}$}{$\dfrac{-7}{3}$}{$\dfrac{-3}{7}$}
\mcq{200}{Solution of the equation 3x+6/5 =0 is:}{$\dfrac{1}{5}$}{$\dfrac{-2}{5}$}{$\dfrac{2}{5}$}{$\dfrac{-1}{5}$}
\mcq{201}{If 2x-7=13 then of x is:}{20}{10}{30}{5}
\mcq{202}{If 5x+70=0 then x=:}{12}{-12}{14}{-14}
\mcq{203}{Linear equation is one variable has:}{One root}{Two roots}{Three roots}{Four roots}
\mcq{204}{The solution of equation 5x+1=0:}{5}{-5}{$\dfrac{-1}{5}$}{$\dfrac{1}{5}$}
\mcq{205}{If 6x+90=0 then x is equal to:}{15}{20}{10}{-15}
\mcq{206}{If 3x-12=0 then x is equal to:}{-4}{6}{12}{4}
\mcq{207}{Given that x+(x+8)=20 the value of x is:}{4}{6}{8}{10}
\mcq{208}{If 6x-3=0 what will be the value of x:}{2}{9}{16}{$\dfrac{1}{2}$}
\mcq{209}{If x+4=12 then value of x is:}{8}{4}{16}{12}
\mcq{210}{In the following, linear equation is:}{$3x^{2}+5x-2=0$}{$\dfrac{2}{x}+\dfrac{3}{y}=5$}{2x+3y=5}{x+2y+3xy=5}
\mcq{211}{The value of x in $\dfrac{20}{7}=\dfrac{80}{x}$ is:}{2}{24}{26}{28}
\mcq{212}{7 times a number is 1085:}{185}{215}{245}{155}
\mcq{213}{If 9 times a number is 180 then the number is:}{10}{20}{30}{40}
\mcq{214}{5x+7=0 is a:}{Non-linear equation}{Quadratic equation}{Cubic equation}{Linear equation}
\mcq{215}{The equation ax+b=0 is:}{Linear}{Quadratic}{Constant}{Non-linear}
\mcq{216}{Degree of quadratic equation is:}{One}{Two}{Three}{Four}
\mcq{217}{The roots of the equations ax$^{2}$+bx+c=0 are real and unequal if:}{$b^{2}=4ac$}{$b^{2}>4ac$}{$b^{2}<4ac$}{$b^{2}=ac$}
\mcq{218}{Product of roots of equation ax$^{2}$+bx+c=0 :}{$\dfrac{a}{b}$}{$\dfrac{-b}{a}$}{$\dfrac{-c}{a}$}{$\dfrac{c}{a}$}
\mcq{219}{Roots are real if:}{$b^{2}>0$}{$a^{2}-b^{2}<0$}{$b^{2}-4ac>0$}{$b^{2}-4ac<0$}
\mcq{220}{The roots of quadratic is equal:}{$b^{2}-4ac<0$}{$b^{2}-4ac<0$}{$b^{2}-4ac=0$}{$b^{2}-4ac\neq0$}
\mcq{221}{The roots of quadratic equation 2x$^{2}$-5x+3=0 are:}{Real and equal}{Real and unequal}{Imaginary}{None of these}
\mcq{222}{A quadratic equation is also called an equation of degree:}{0}{1}{2}{3}
\mcq{223}{A quadratic equation at the most:}{One root}{Two roots}{Three roots}{Four roots}
\mcq{224}{A second degree equation is called:}{Linear equation}{Quadratic equation}{Cubic equation}{Radical equation}
\mcq{225}{$x=\dfrac{\sqrt{b^{2}-4ac}}{2a}$ is quadratic fromula for quadratic equation:}{$ax^{2}-bx-c=0$}{$ax^{2}+bx+c=0$}{$ax^{2}+bx-c=0$}{$ax^{2}-bx+c=0$}
\mcq{226}{Standard formula of quadratic equation is:}{ax+b=0}{$ax^{2}+bx=0$}{$ax^{2}=0$}{$ax^{2}+bx+c=0$}
\mcq{227}{In quadratic equation the variable has degree ......:}{0}{1}{2}{None}
\mcq{228}{Number of ways to solve a quadratic equation are:}{Two}{Three}{Four}{None of these}
\mcq{229}{Degree of quadratic equation is:}{One}{Two}{Three}{Four}
\mcq{230}{The degree of the equation x$^{2}$-2x+1=0 is:}{4}{3}{2}{1}
\mcq{231}{If ax$^{2}$+bx+c=0, a is not equal to 0 then:}{$x=\dfrac{-b\pm\sqrt{4ac-b^{2}}}{2}$}{$x=\dfrac{-b\pm\sqrt{b^{2}-4ac}}{2}$}{$x=\dfrac{-b\pm\sqrt{b^{2}-4ac}}{2a}$}{$x=\dfrac{-b\pm\sqrt{b^{2}+4ac}}{2a}$}
\mcq{232}{Formula $x=\dfrac{-b\pm\sqrt{b^{2}-4ac}}{2a}$ is called:}{Distance formula}{Quadratic formula}{Interest formula}{Slope formula}
\mcq{233}{Maximum terms in a quadratic equation are:}{2}{4}{6}{3}
\mcq{234}{The factors of x$^{2}$+3x+2=0 are:}{\{x-2\}\{x-3\}}{\{x-2\}\{x+3\}}{\{x+2\}\{x-1\}}{\{x+2\}\{x+1\}}
\mcq{235}{The solution set of $\sqrt{x}+5=4$ is:}{\{1\}}{\{-1\}}{\{ \}}{\{+1\}}
\mcq{236}{The root of x$^{2}$+2x+1=0:}{-1}{2}{-2}{1}
\mcq{237}{Solution of the equation x$^{2}$-6x+9=0:}{1}{3}{5}{9}
\mcq{238}{2x$^{2}$+3x+1=0 is:}{Linear equation}{Cubic equation}{Quadratic equation}{None of these}
\mcq{239}{The solution set of equation x$^{2}$+13x-90=0 is:}{\{18,5\}}{\{9,5\}}{\{-18,5\}}{\{9,4\}}
\mcq{240}{The roots of x$^{2}$-3x+2=0:}{\{0,1\}}{\{1,1\}}{\{1,2\}}{\{2,2\}}
\mcq{241}{The equation 2x$^{2}$-3x+7=0 is of degree:}{1}{3}{2}{4}
\mcq{242}{Solution set of equation x$^{2}$-4=0 is:}{\{2\}}{\{-2,2\}}{\{-5\}}{25}
\mcq{243}{Solution set of x$^{2}$-25=0 is:}{\{-7\}}{\{7\}}{\{-7,7\}}{\{49\}}
\mcq{244}{The solution set of 4x$^{2}$+5x+1=0 is:}{$\left\{1,\dfrac{1}{4}\right\}$}{\{1,4\}}{\{-1,-4\}}{$\left\{-1,\dfrac{1}{4}\right\}$}
\mcq{245}{Solution of the quadratic equation x$^{2}$-6x+8=0 are:}{\{-6,8\}}{\{6,-8\}}{\{2,4\}}{\{-2.-4\}}
\mcq{246}{Solution set of x$^{2}$-9=0 is:}{\{3\}}{\{-3\}}{9}{\{-3,3\}}
\mcq{247}{In quadratic equation the varaible has degree:}{1}{2}{More than 2}{Less than 2}
\mcq{248}{Discriminant in quadratic formula is:}{b$^{2}$-4ac}{b$^{2}$}{-4ac}{$\dfrac{b^{2}-4ac}{2a}$}
\mcq{249}{The roots of quadratic equation 2x$^{2}$-5x3=0 are:}{Real and equal}{Real and unequal}{Imaginary}{None of these}
\mcq{250}{An equation of the form Ia$^{x}$+ma$^{-x}$+n=0 is called ...... equation:}{Linear}{Quadratic}{Reciprocal}{Exponential}
\mcq{251}{Root of the equation that do not satisfy the equation are called:}{Internal root}{Extraneous root}{Outer root}{Real root}
\mcq{252}{The equation of the type 3$^{2x}$+9=10.3x is:}{Irrational equation}{Reciprocal equation}{Exponential equation}{None of these}
\mcq{253}{The point ot remember for reciprocal equation is that the coefficient from the beginning and end are:}{equal}{Not equal}{May or may not equations}{None of these}
\mcq{254}{System of simultaneous equation is solved by:}{factorization}{Subtraction or addition}{Substitution}{Both b and c}
\mcq{255}{Sum and product of two number is 25,100 when:}{15,10}{20,5}{30,-5}{25,4}
\mcq{256}{The number of order pair of solution set of x+y=11 and x-y=7 are:}{One}{Two}{Three}{Four}
\mcq{257}{Solution set of x+y=2 and 2x-y=7 are:}{\{-3,1\}}{\{-3,-1\}}{\{3,-1\}}{None of these}
\mcq{258}{If x+y=6, x-y=10 then the value of x and y are:}{x=6,y=10}{n=4, y=6}{x=6,y=4}{x=8,y=-2}
\mcq{259}{Solution set of equation x+y=0, x-y=0 is:}{\{(0,0)\}}{\{(1,-1)\}}{\{(-1,1)\}}{\{(-1,-1)\}}
\mcq{260}{The solution set of x+y=12, x-y=-2 is:}{x=1,y=2}{x=4,y=3}{x=5,y=4}{x=5,y=7}
\mcq{261}{Solution of simultaneous equation x-y=4 and x+y=8 is:}{x=1,y=2}{x=2,y=3}{x=6,y=1}{x=6,y=2}
\mcq{262}{Given x=y and 2x+y=3 then solution set is:}{\{(1,1)\}}{\{(1,2)\}}{\{(2,1)\}}{\{(2,2)\}}
\mcq{263}{Solution set of x+y=2 and x-y=0 is:}{\{(0,0)\}}{\{(1,0)\}}{\{(0,1)\}}{\{(1,1)\}}
\mcq{264}{Solve y-x=2 amd y+x=2 for y:}{4}{2}{0}{-2}
\mcq{265}{Solution of the simultaneous equations x+y=1 ; x-y= 1 is:}{\{1,1\}}{\{1,-1\}}{\{-1,1\}}{\{1,0\}}
\mcq{266}{Solution of equation x+y=5 and x-y=1 is:}{\{3,2\}}{\{2,3\}}{\{-3,-2\}}{\{-2,-3\}}
\mcq{267}{The solution of 2x+y=6 , x-y=2 is:}{$\left[\dfrac{8}{3},\dfrac{2}{3}\right]$}{$\left[\dfrac{3}{8},\dfrac{3}{2}\right]$}{$\left[\dfrac{-3}{8},\dfrac{3}{2}\right]$}{$\left[\dfrac{-8}{3},\dfrac{-2}{3}\right]$}
\mcq{268}{Solution set of x+y=12, x-y=8 is:}{\{(10,2)\}}{\{(4,2)\}}{\{(8,1)\}}{\{(4,6)\}}
\mcq{269}{If 2x+y=3 , 8x+5y=6 then:}{x=0, y=3}{x=4,y=1}{x=2, y=0}{x=1, y=1}
\mcq{270}{The term function was introduce by:}{Newton}{Cauchy}{Leibniz}{Lagrange}
\mcq{271}{The concept of function was introduced in:}{17th century}{18th century}{19th century}{20th century}
\mcq{272}{The term g(x), h(x), p(x) etc are used to indicate:}{Variable}{Constants}{Function rule}{None of these}
\mcq{273}{Domian is the set of all possible value of:}{Input}{Output}{Negative value}{Positive value}
\mcq{274}{If f(x) = c then f is called the function:}{Implicit}{Constant}{Quadratic}{Explicit}
\mcq{275}{The function $f(x)=\dfrac{1}{x}$ is not defined as:}{1}{-1}{0}{2}
\mcq{276}{The set of all values of Y is called if f(x)=y:}{Domain}{Range}{Both a and}{Variable}
\mcq{277}{If f(x,y)=0 the f is called ...... function:}{Constant}{Inplict}{Explict}{Onto}
\mcq{278}{If y=f(x) then f is called:}{Implicit function}{Explicit funtion}{Constant function}{None of these}
\mcq{279}{If f(-x)= f(x) then f is called:}{Exponential}{Odd function}{Even function}{Linear function}
\mcq{280}{The function f(x) is called odd if:}{f\{-x\}= fx}{f\{-x\}=fx}{f\{-x\}=0}{None}
\mcq{281}{The domain of f(x)=\{(a,1),(b,1),(c,1)\} is:}{\{a,b,c\}}{\{a\}}{\{b\}}{\{c\}}
\mcq{282}{The domain of the function $f(x)=\dfrac{1}{2x-3}$ is:}{$R-\left\{\dfrac{3}{2}\right\}$}{$R-\left\{\dfrac{-3}{2}\right\}$}{$R+\left\{\dfrac{3}{2}\right\}$}{$R+\left\{\dfrac{-3}{2}\right\}$}
\mcq{283}{The domain of $f(x)=\dfrac{1}{x+5}$ is:}{R+\{5\}}{R-\{-5\}}{R-\{5\}}{None}
\mcq{284}{If the set \{(2,4),(3,7),(5,8)\} show the raltion x=g(y) then set \{4,7,8\} represent:}{Domain}{Range}{Relation}{None}
\mcq{285}{A cubic function has degree:}{2}{1}{0}{3}
\mcq{286}{The domain of the function $f(x)=\dfrac{1}{x-3}$ is:}{$R-\{x|x=3\}$}{$R-\{x|x\neq\}$ [text cut off in the source]}{3}{R}
\mcq{287}{If F(x)=2-x$^{2}$the F(-8) equals:}{61}{63}{66}{-62}
\mcq{288}{$f(x)=\dfrac{x}{x+2}$ then f(-1) is:}{0}{1}{-1}{$\dfrac{1}{2}$}
\mcq{289}{The function value $h\left(\dfrac{1}{16}\right)$ of function $h(u)=\dfrac{1}{\sqrt{u}}$ is:}{$\dfrac{1}{4}$}{4}{16}{$-\dfrac{1}{4}$}
\mcq{290}{If $f(x)=\dfrac{x^{2}}{4}+6$ the f(-2) is equal is:}{7}{3}{8}{-4}
\mcq{291}{If $f(x)=\sqrt{2x+1}$ , then f(760) is:}{29}{39}{49}{None of these}
\mcq{292}{The function $g(t)=\dfrac{1}{20}$ is:}{Linear function}{Quardratic function}{Constant function}{None of these}
\mcq{293}{If $f(x)=\dfrac{1}{x-4}$ then f(3) is:}{-1}{Undefined}{Zero}{1}
\mcq{294}{The value of fucntion f(x)=5x-30 at x=6 is:}{0}{-1}{1}{8}
\mcq{295}{If f(x)=x$^{2}$-4x+1 then f(-2) is equal to:}{13}{3}{12}{-5}
\mcq{296}{If f(x) = 3x-6 then f(0) is:}{6}{-6}{3}{-3}
\mcq{297}{If f(x)= x$^{2}$-2x+3 then f(-1) is equal to:}{Zero}{6}{4}{2}
\mcq{298}{If f(x) = x$^{2}$-x+1 then f(2) is equal to:}{3}{Zero}{5}{-3}
\mcq{299}{If the set \{(2,4),(3,7),(5,8)\} shows a relation y=f(x) then the set \{2,3,5\} represents:}{Range}{Domain}{Function}{Prime numbers}
\mcq{300}{If f(x)= x$^{2}$-4 then f(-2) is equal to:}{8}{6}{2}{Zero}
\mcq{301}{If f(x) = 2x+26 then f(2)):}{63}{64}{68}{66}
\mcq{302}{Value of f(x)=x$^{2}$+2x+1 at x=2 is:}{4}{6}{8}{9}
\mcq{303}{If f(x)= x$^{3}$-1 then f(-1) is equal to:}{Zero}{2}{-2}{-3}
\mcq{304}{The value of the function f(x)=6x-3 at x=2 is:}{8}{9}{7}{None of these}
\mcq{305}{If f(x)=9x$^{2}$+3x+5 then f(-2)=:}{20}{25}{30}{35}
\mcq{306}{y=f(x) then x is called:}{Dependent variable}{Independent variable}{Both A and B}{None of these}
\mcq{307}{If f(-x)=-f(x) then f is called:}{Even function}{Odd function}{Constant function}{None of these}
\mcq{308}{If given function is g(u)=u$^{2}$+u, find the value of at x=2x:}{6}{3c}{4c$^{2}$+2c}{None of these}
\mcq{309}{If f(x)=2x-10, if f(x) = 1 then x is:}{0}{2}{10}{5}
\mcq{310}{If $f(x)=\dfrac{x+5}{x-5}$ , then f(0) is:}{1}{10}{0}{-1}
\mcq{311}{The equation x=0 represents:}{X=axis}{Y-axis}{Origin}{Both and b}
\mcq{312}{For the point P(3,6):}{3 is abscissa}{6 is abscissa}{3 is ordinate}{Both 3 and 6 are ordii}
\mcq{313}{In which quadrant the point (-5,-3) lies:}{I}{II}{III}{IV}
\mcq{314}{The point (-2,6) lies in the quadrant:}{1st}{IInd}{IIIrd}{IVth}
\mcq{315}{In XY-plane, the co-ordinates of origin are:}{\{X,0\}}{\{0,Y\}}{X,Y}{\{0,0\}}
\mcq{316}{The point (-3,-4) lies in quardrant:}{I}{II}{III}{IV}
\mcq{317}{The point (-3,4) lie in the quadrant:}{I}{II}{III}{IV}
\mcq{318}{Fix quantities are known as:}{Constant}{Variable}{Abscissa}{Ordinate}
\mcq{319}{A point(-x,y) lies in:}{1st quadrant}{2nd quadrant}{3rd quadrant}{4th quadrant}
\mcq{320}{The slope of y-5x+3=0 is:}{5}{$\dfrac{5}{3}$}{3}{1}
\mcq{321}{Y-intercept of the line $y=\dfrac{1}{2}(6x+7)$ is:}{7}{5}{$\dfrac{7}{2}$}{6}
\mcq{322}{X-intercept of 2x-3y=6 is:}{3}{-2}{6}{12}
\mcq{323}{Graph of the linear equation is:}{Para bola}{Straight line}{Circle}{Impossible}
\mcq{324}{The x-coordinate of a point is called:}{Ordinate}{Abscissa}{Origin}{Quadrant}
\mcq{325}{The origin is:}{\{0,x\}}{\{y,0\}}{\{0,0\}}{\{x,y\}}
\mcq{326}{The point (-2,6) lies in quadrant:}{1st}{2nd}{3rd}{4th}
\mcq{327}{The y-coordinate of any point is:}{Abscissa}{Ordinate}{x-intercept}{Origin}
\mcq{328}{For 1st quadrant, the values of x and y must:}{Both negative}{Both positive}{x positive, y negative}{x negative, y positive}
\mcq{329}{Graph of a linear function is always:}{Straight line}{Parabola}{Curve}{None of these}
\mcq{330}{Graph of a quadratic function is:}{Straight line}{Parabola}{Curve}{None of these}
\mcq{331}{A rectangular array of number enclosed by a pair of brackets is called:}{Set}{Schedule}{Matrix}{Scalar}
\mcq{332}{The order of the matix A\{1 2 5 8\} is:}{4 x 1}{4 x 4}{1 x 4}{1 x 1}
\mcq{333}{A matrix \{ 1 2 3\} is called:}{Column}{Square}{Zero}{Row}
\mcq{334}{The horizintal line of a number in a matrix are called:}{Columns}{Rows}{Vectors}{Scalars}
\mcq{335}{The vertical line of numbers in a matrix is called:}{Columns}{Rows}{Vectors}{Scalars}
\mcq{336}{The order matrix $\begin{bmatrix}4\\1\\3\end{bmatrix}$ is:}{1 x 3}{3 x 1}{4 x 3}{1 x 1}
\mcq{337}{A matrix in which each element is zero is called:}{Identity matrix}{Scalar matrix}{Null matrix}{Non- singular}
\mcq{338}{The order of the matrix $\begin{bmatrix}1&4\\2&5\\3&6\end{bmatrix}$ is:}{3 x 2}{2 x 3}{2 x 2}{3 x 3}
\mcq{339}{$\begin{bmatrix}K&0\\0&K\end{bmatrix}$ is a ...... matrix:}{Zero}{Non-diagonal}{Identity}{Scalar}
\mcq{340}{If A is order m x n and B is of order n x m matrix then AB will be of order:}{m x n}{n x m}{m x m}{n x n}
\mcq{341}{The order of matrix $\begin{bmatrix}2\\5\end{bmatrix}$ is:}{2 x 2}{2 x 1}{1 x 2}{1 x 1}
\mcq{342}{$\begin{bmatrix}K&O&O\\O&K&O\\O&O&K\end{bmatrix}$ is ...... matrix:}{Diagonal}{Row}{Scalar}{None}
\mcq{343}{$\begin{bmatrix}3&2\\1&0\\1&3\end{bmatrix}$ so order of A is:}{3 x 2}{2 x 3}{2 x 2}{3 x 3}
\mcq{344}{The matrix \{ a b c\} is:}{Diagonal matrix}{Scalar}{Row matrix}{Identity}
\mcq{345}{The matrix $\begin{bmatrix}2&0\\0&3\end{bmatrix}$ is:}{Scalar}{Diagonal}{Identity}{Null}
\mcq{346}{The order of matrix $\begin{bmatrix}1&3\\2&4\end{bmatrix}$ is:}{2 x 2}{2 x 1}{1 x 2}{None}
\mcq{347}{A matrix having x rows and n column is called square matrix if:}{m $>$ n}{m$<$ n}{m = n}{None}
\mcq{348}{The order of the matrix [ 1 2 8 3] is:}{4 x 1}{1 x 4}{4 x 4}{3 x 4}
\mcq{349}{Order of the matrix [ 1 2 8 3] is:}{3 x 1}{3 x 3}{1 x 3}{1 x 2}
\mcq{350}{matrix having single row is called:}{Row matrix}{Singular matrix}{Squarematrix}{Null matrix}
\mcq{351}{Order of the matrix having m row and n columns is:}{m - n}{m x n}{n x m}{$\dfrac{m}{n}$}
\mcq{352}{A matrix in which the number of rows are same as that of columns is called:}{Scalar matrix}{Rectangular matrix}{Square matrix}{Column matrix}
\mcq{353}{$\begin{bmatrix}1&0\\0&1\end{bmatrix}$ is called:}{Rectangular matrix}{Diagonal matrix}{Identity matrix}{Singular matrix}
\mcq{354}{In order of matrix A is 2 x 2 and order of matrix B is 2 x 1 then order of (AB)$^{t}$ will be:}{2 x 1}{1 x 2}{2 x 2}{1 x 1}
\mcq{355}{Order of matrix $\begin{bmatrix}3&1&4\\2&0&-1\end{bmatrix}$ is:}{3 x 2}{3 x 3}{2 x 2}{2 x 3}
\mcq{356}{[2 -1 3 4] is a ...... matrix:}{Row}{Column}{Null}{None}
\mcq{357}{$\begin{bmatrix}1&0\\0&1\end{bmatrix}$ is identity matrix of order:}{0 x 0}{1 x 1}{2 x 2}{None}
\mcq{358}{Matrix has only one column is called:}{Square matrix}{Singular}{Column Matrix}{Row matrix}
\mcq{359}{The order of matrix is: [matrix missing in the source]}{1 x 2}{2 x 1}{2 x 2}{2 x 3}
\mcq{360}{The order of $\begin{bmatrix}1\\2\\8\\3\end{bmatrix}$ is:}{4 x 1}{1 x 4}{1 x 1}{4 x 4}
\mcq{361}{The order of transpose of $\begin{bmatrix}1&2\\3&4\\5&6\end{bmatrix}$ is:}{2 x 2}{2 x 3}{3 x 2}{3 x 3}
\mcq{362}{A matrix which has one column is called:}{Row matrix}{Column matrix}{Singular matrix}{Vector matrix}
\mcq{363}{Which one is not true in general:}{A+B=B+A}{AB=BA}{A+}{A}
\mcq{364}{If 3A + $\begin{bmatrix}2&5\\4&6\end{bmatrix}$ =0 then A=......:}{$\begin{bmatrix}\frac{2}{3}&\frac{5}{3}\\\frac{4}{3}&2\end{bmatrix}$}{$\begin{bmatrix}-1&-2\\1&3\end{bmatrix}$}{$\begin{bmatrix}\frac{2}{3}&\frac{5}{3}\\4&6\end{bmatrix}$}{$\begin{bmatrix}\frac{-2}{3}&\frac{-5}{3}\\\frac{-4}{3}&-2\end{bmatrix}$}
\mcq{365}{If $A=$ $\begin{bmatrix}3&0\\6&0\end{bmatrix}$, then $|A|$=......:}{3}{0}{6}{18}
\mcq{366}{If A=[0] then A$^{-1}$ is:}{0}{1}{Not}{Not possible}
\mcq{367}{The matrix A will by skew-symmetric if:}{A+A$^{t}$}{A=A$^{-1}$}{A$^{t}$=-A}{A=A$^{t}$}
\mcq{368}{(AB)$^{t}$=......:}{AB}{BA}{A$^{t}$B$^{T}$}{B$^{t}$A$^{t}$}
\mcq{369}{If A = $\begin{bmatrix}2&3\\1&5\end{bmatrix}$ and B = $\begin{bmatrix}1&7\\6&4\end{bmatrix}$ then A+B=:}{$\begin{bmatrix}3&10\\7&9\end{bmatrix}$}{$\begin{bmatrix}9&4\\5&11\end{bmatrix}$}{$\begin{bmatrix}4&5\\9&11\end{bmatrix}$}{$\begin{bmatrix}10&3\\9&7\end{bmatrix}$}
\mcq{370}{If A = $\begin{bmatrix}-2&6\\4&7\end{bmatrix}$, then A$^{t}$ is:}{$\begin{bmatrix}-2&7\\6&4\end{bmatrix}$}{$\begin{bmatrix}7&6\\4&-2\end{bmatrix}$}{$\begin{bmatrix}-2&4\\6&7\end{bmatrix}$}{$\begin{bmatrix}-2&7\\6&4\end{bmatrix}$}
\mcq{371}{If A = $\begin{bmatrix}1&0\\0&1\end{bmatrix}$ then $|A|$ is:}{1}{0}{2}{3}
\mcq{372}{If two rows in a squar matrix are identical then its determinant is:}{-1}{-2}{2}{0}
\mcq{373}{If order of A is m x n then order of A$^{t}$ is:}{m x n}{n x m}{m x m}{n x m}
\mcq{374}{A square matrix is symmetric if A$^{t}$ equals:}{-A}{A$^{2}$}{A}{-A$^{t}$}
\mcq{375}{(A+B)$^{t}$ :}{Bt-A$^{t}$}{A$^{t}$+b$^{t}$}{A$^{t}$B$^{t}$}{A$^{t}$-B$^{t}$}
\mcq{376}{For any square matrix A, $|A^{t}|$ will be equal to:}{A$^{-1}$}{Adj}{$|A|$}{A$^{t}$}
\mcq{377}{The transpose of a matrix is denoted by:}{A}{A$^{t}$}{A$^{-1}$}{$|A|$}
\mcq{378}{For a symmetric matrix if A$^{t}$ is equal to:}{-A}{0}{0}{A}
\mcq{379}{If A = $\begin{bmatrix}2&9\\3&8\end{bmatrix}$ then $|A|$ is:}{11}{-11}{42}{43}
\mcq{380}{if A = $\begin{bmatrix}12&8\\2&4\end{bmatrix}$ then $\begin{bmatrix}6&4\\1&2\end{bmatrix}$ is equal to:}{12}{$\dfrac{1}{2}A$}{8}{32}
\mcq{381}{The transpose of a row matrix is always a:}{Column matrix}{Row matrix}{Square matrix}{Null matrix}
\mcq{382}{Any square matrix A is a symmetric matrix if:}{A=-A}{A=A$^{t}$}{A=-A$^{t}$}{A=A$^{-t}$}
\mcq{383}{If A and B are two matrixs then (AB)t:}{ABB}{BA}{B$^{t}$A$^{T}$}{A$^{t}$B$^{t}$}
\mcq{384}{The transpose of row matrix is a:}{Row matrix}{Column matrix}{Unit matrix}{None of these}
\mcq{385}{A square matrix having all entries above the principle digonal as zero is called:}{Null matrix}{Upper triangular matrix}{Lower triangularmatrix}{Diagonal matrix}
\mcq{386}{The order of a matrix $\begin{bmatrix}1&3\\8&5\end{bmatrix}$ is:}{3 x 2}{2 x 2}{2 x 3}{1 x 3}
\mcq{387}{The matrix $\begin{bmatrix}2&0\\0&3\end{bmatrix}$ is:}{Scalar matrix}{Identity matrix}{Diagonal matrix}{Null matrix}
\mcq{388}{The matrix $\begin{bmatrix}4&4\\6&6\end{bmatrix}$ is:}{Null matrix}{Singular matrix}{Diagonal matrix}{Null matrix}
\mcq{389}{If $A=\begin{bmatrix}a_{11}&a_{12}\\a_{21}&a_{22}\end{bmatrix}$, then the matrix, $\begin{bmatrix}a_{22}&-a_{12}\\-a_{21}&a_{11}\end{bmatrix}$ is called:}{Adjoint of A}{Transpose of A}{Inverse of A}{Square of A}
\mcq{390}{If X + $\begin{bmatrix}3\\4\end{bmatrix}$ = $\begin{bmatrix}-4\\3\end{bmatrix}$, then matrix will be:}{$\begin{bmatrix}0\\0\end{bmatrix}$}{$\begin{bmatrix}-4\\3\end{bmatrix}$}{$\begin{bmatrix}-4\\-3\end{bmatrix}$}{$\begin{bmatrix}-7\\-1\end{bmatrix}$}
\mcq{391}{A matrix A is called summetric if:}{A$^{t}$=-A}{A$^{t}$=A}{A$^{t}$=|A|}{None of these}
\mcq{392}{Order of the matrix $\begin{bmatrix}1&3&-1\\6&1&8\end{bmatrix}$ is:}{3 x 3}{3 x 2}{2 x 3}{1 x 4}
\mcq{393}{[formula missing in the source] is called:}{Transpose of A}{Inverse of A}{Minor of A}{Adjoint of A}
\mcq{394}{A square matrix A is said to be singular if ......:}{$|A|\neq0$}{|A|=0}{|A|<0}{|A|>0}
\mcq{395}{If A is matrix or order 3 x 4 and B is a matrix of 4 x 5, then AB matrix will be or order ......:}{3 x 3}{3 x 4}{3 x 5}{4 x 4}
\mcq{396}{If A is order m x n and B is of order n x m matrices, then AB will be of order:}{m x n}{n x m}{m x m}{n x n}
\mcq{397}{If the matrix $\begin{vmatrix}K&4\\3&2\end{vmatrix}$ is singular then:}{K=6}{K=4}{K=3}{K=2}
\mcq{398}{Adjoint of $\begin{bmatrix}2&-3\\-1&4\end{bmatrix}$ is:}{$\begin{bmatrix}2&-1\\-3&-4\end{bmatrix}$}{$\begin{bmatrix}-2&-1\\-3&4\end{bmatrix}$}{$\begin{bmatrix}4&3\\1&2\end{bmatrix}$}{None of these}
\mcq{399}{If A= $\begin{bmatrix}a&b\\c&d\end{bmatrix}$ then adj (A) is:}{$\begin{bmatrix}d&-b\\-c&a\end{bmatrix}$}{$\begin{bmatrix}d&c\\b&a\end{bmatrix}$}{$\begin{bmatrix}a&-b\\-c&d\end{bmatrix}$}{$\begin{bmatrix}-a&c\\b&-d\end{bmatrix}$}
\mcq{400}{If a = $\begin{bmatrix}1&1\\2&3\end{bmatrix}$ the |A| is:}{0}{1}{2}{3}
\mcq{401}{Square matrix A is called singular if:}{|A|=1}{$|A|\neq0$}{|A|=0}{$|A|\neq1$}
\mcq{402}{2 $\begin{vmatrix}1&2\\3&4\end{vmatrix}$ =:}{20}{-8}{-4}{40}
\mcq{403}{The matrix $\begin{bmatrix}3x&4\\3&-2\end{bmatrix}$ will be singular for value of x:}{0}{-2}{2}{3}
\mcq{404}{If A= $\begin{bmatrix}2&-1\\4&3\end{bmatrix}$ then |A|=:}{6}{10}{-4}{8}
\mcq{405}{A= $\begin{bmatrix}2&-2\\4&-4\end{bmatrix}$ then |A|=:}{8}{16}{0}{32}
\mcq{406}{The determinant of an identity matrix is:}{0}{1}{-1}{2}
\mcq{407}{The matrix $\begin{bmatrix}2&2\\3&3\end{bmatrix}$ is:}{Singular}{Non-singular}{Scalar}{Rectangular}
\mcq{408}{If A= $\begin{bmatrix}X&1\\1&1\end{bmatrix}$,|A|=0 then x=:}{0}{1}{2}{3}
\mcq{409}{$\begin{vmatrix}X&4\\5&10\end{vmatrix}$=0 then X=:}{2}{4}{6}{8}
\mcq{410}{If AX=B then X is:}{BA}{A$^{-1}$B}{AB}{AB$^{-1}$}
\mcq{411}{The additive inverse of matrix A is:}{A}{-A}{A$^{-1}$}{$\dfrac{1}{A}$}
\mcq{412}{$\begin{bmatrix}3&5\\1&4\end{bmatrix}$ matrix is:}{None-singular}{Singular}{Identity}{All}
\mcq{413}{The matrix $\begin{bmatrix}4&4\\6&6\end{bmatrix}$ is:}{Null matrix}{Singular matrix}{Diagonal}{Unit matrix}
\mcq{414}{If $\begin{vmatrix}3&1\\2&x\end{vmatrix}$=0; then x=:}{$\dfrac{1}{3}$}{$\dfrac{1}{4}$}{$\dfrac{2}{3}$}{$\dfrac{3}{2}$}
\mcq{415}{Inverse of square matix A is not possible if:}{|A|=1}{|A|=2}{|A|=0}{$|A|\neq0$}
\mcq{416}{The determinant of matrix $\begin{bmatrix}2&1\\-1&0\end{bmatrix}$ is equal to:}{2}{-1}{2}{1}
\mcq{417}{If matrix A= $\begin{bmatrix}2&1\\2&1\end{bmatrix}$, the |A| is:}{4}{0}{-4}{-8}
\mcq{418}{A square matrix A is called non-singular if:}{|A|=0}{$|A|\neq0$}{|A|=|A$^{t}$|}{|A|=-|A|}
\mcq{419}{If B= $\begin{bmatrix}-2&-4\\-3&4\end{bmatrix}$ then |B| is equal to:}{-20}{20}{10}{-10}
\mcq{420}{Inverse of a square matrix A is only possible when:}{|A|=0}{$|A|\neq0$}{|A|=1}{|A|=-1}
\mcq{421}{If matrixes A and B are non singular, then (AB)$^{-1}$=1:}{A$^{-1}$B$^{-1}$}{AB}{B$^{-1}$A$^{-1}$}{AB}
\mcq{422}{IF $\begin{vmatrix}3&6\\3&x\end{vmatrix}$=0 then x is:}{6}{18}{-6}{9}
\mcq{423}{If A= $\begin{bmatrix}3&2\\1&-1\end{bmatrix}$, then |A| is equal to:}{-5}{5}{1}{6}
\mcq{424}{If matrix $\begin{bmatrix}x&2\\1&2\end{bmatrix}$ is singular then x is equal to:}{Xero}{1}{-1}{2}
\mcq{425}{If A = $\begin{bmatrix}-2&3\\4&2\end{bmatrix}$ then |A| is equal to:}{-2}{Zero}{-16}{16}
\mcq{426}{If A = $\begin{bmatrix}5&-2\\3&1\end{bmatrix}$ then |A| is:}{11}{-1}{13}{17}
\mcq{427}{If A = $\begin{bmatrix}1&2\\3&4\end{bmatrix}$ then |A| is:}{-2}{6}{4}{None of these}
\mcq{428}{$\begin{bmatrix}1&1\\1&1\end{bmatrix}$ matrix is:}{Null}{Singular}{Identity}{Diagonal}
\mcq{429}{If A = $\begin{bmatrix}2&4\\6&x\end{bmatrix}$ is singular matrix then x is:}{12}{6}{8}{24}
\mcq{430}{If $\begin{bmatrix}5x&8\\5&2\end{bmatrix}$=0 then x is equal to:}{1}{2}{-4}{4}
\mcq{431}{The value of the determinant of every identity matrix is always:}{0}{1}{2}{3}
\mcq{432}{An inverse matrix is always a:}{Singular martix}{Row matrix}{Non-singular matrix}{Column matrix}
\mcq{433}{Product of a matrix and its inverse is always:}{Identity matrix}{A square matrix}{A row matrix}{A column matrix}
\mcq{434}{B is not equal to zero, then B is called:}{Singular matrix}{Nor-singular matrix}{Matrix only}{Null matrix}
\mcq{435}{If AX=B then X is:}{BA$^{-1}$}{AB}{$\dfrac{B}{A}$}{A$^{-1}$B}
\mcq{436}{The base of binary system is:}{4}{8}{10}{2}
\mcq{437}{The binary number system is based on ......:}{2}{4}{8}{10}
\mcq{438}{The binary number system to based on digits:}{[1,2]}{[0,2]}{[0,1]}{[3,4]}
\mcq{439}{Decimal number system is based on digits:}{Two}{Four}{Eight}{Ten}
\mcq{440}{Decimal number system is based on the digits:}{0 to 2}{0 to 3}{0 to 5}{0 to 9}
\mcq{441}{The number system used in our every day life is:}{Binary}{Decimal}{Octal}{Hexadecimal}
\mcq{442}{Decimal number system is based on the digits:}{0 to 9}{0 to 1}{1 to 10}{0 to 15}
\mcq{443}{A decimal number system is also called:}{Binary}{Denary}{Greek}{Octal}
\mcq{444}{(9)10 in binary system is:}{[1010]$_{2}$}{[1001]$_{2}$}{[1100]$_{2}$}{[1000]$_{2}$}
\mcq{445}{Arabian system is also called:}{Decimal}{Quantal}{Binary}{Octal}
\mcq{446}{Arabian number system is based on digits:}{2}{4}{8}{10}
\mcq{447}{(1000)$_{2}$ in decimal number system is:}{8}{10}{16}{12}
\mcq{448}{In computer, the number system used is:}{Octal}{Decimal}{Hexa}{Binary}
\mcq{449}{In the binary system 2 equals to:}{10$_{2}$}{11$_{2}$}{101$_{2}$}{110$_{2}$}
\mcq{450}{In computer 0 \& 1 respectively is used for:}{On only}{Off only}{On \& off}{Off \& on}
\mcq{451}{The number 4 in binary number system is:}{101$_{2}$}{100$_{2}$}{111$_{2}$}{1010}
\mcq{452}{In the binary number system 5=\underline{\hspace{3em}}:}{100$_{2}$}{110$_{2}$}{11$_{2}$}{101$_{2}$}
\mcq{453}{In binary system 8 is=\underline{\hspace{3em}}:}{100$_{2}$}{1000$_{2}$}{10000}{None}
\mcq{454}{In the binary system 3 is equivalent to:}{11$_{2}$}{10$_{2}$}{01$_{2}$}{101$_{2}$}
\mcq{455}{7 in binary number system is:}{11$_{2}$}{101$_{2}$}{110$_{2}$}{111$_{2}$}
\mcq{456}{Convert (6)10 into binary number:}{1001$_{2}$}{1000$_{2}$}{111$_{2}$}{110$_{2}$}
\mcq{457}{Convert 37 into binary number is:}{100111$_{2}$}{111101$_{2}$}{100101$_{2}$}{111000$_{2}$}
\mcq{458}{(1011)$_{2}$ is decimal system is:}{11}{15}{19}{21}
\mcq{459}{The decimal number of (1001)2 is:}{7}{9}{11}{13}
\mcq{460}{Decimal for of (101)$_{2}$ is:}{2}{3}{5}{1}
\mcq{461}{(111)$_{2}$ in decimal system is:}{3}{5}{7}{9}
\mcq{462}{The equivalent number of (10)$_{2}$ in decimal system is:}{4}{3}{5}{2}
\mcq{463}{The binary number (11)$_{2}$ is equal to decimal number:}{22}{5}{7}{3}
\mcq{464}{In decimal system, (110)$_{2}$ is equal to:}{8}{2}{4}{6}
\mcq{465}{In decimal system the number (1010)$_{2}$ is:}{12}{21}{08}{10}
\mcq{466}{The binary number 11 is in decimal system is equivalent to:}{1}{2}{3}{4}
\mcq{467}{In decimal system (10001)$_{2}$ is equivalent to the number:}{15}{16}{17}{18}
\mcq{468}{(101)$_{2}$ equal to decimal:}{101}{202}{10}{5}
\mcq{469}{(11)10 = ......:}{1101$_{2}$}{1011$_{2}$}{1010$_{2}$}{1111$_{2}$}
\mcq{470}{(10110)$_{2}$ in decimal number is:}{20}{24}{22}{26}
\mcq{471}{(11)$_{2}$+(10)$_{2}$ is equal to:}{100$_{2}$}{110$_{2}$}{101$_{2}$}{1010$_{2}$}
\mcq{472}{To convert a decimal number into binary number we use the method of:}{Addition}{Division}{Multiplication}{Subtraction}
\mcq{473}{The decimal number 43 comprises:}{4 units and 3 tens}{3 tens and no unit}{0 unit and 4 tens}{3 unit and 4 tens}
\mcq{474}{The decimal number 23 in simplest form is:}{3[10]$^{0}$+2[10]$^{1}$}{2[10]$^{0}$+3[10]$^{1}$}{2[10]$^{0}$+3[10]$^{2}$}{3[10]+2[10]$^{2}$}
\mcq{475}{Octal number system is based on:}{Two digits}{Ten digits}{Eight digits}{Sixteen digits}
\mcq{476}{Hexadecimal number system is based on:}{Two digits}{Ten digits}{Eight digits}{Sixteen digits}
\mcq{477}{(1001)$^{2}$ in decimal system is:}{9}{7}{13}{11}
\mcq{478}{Conversion of 37 into binary number is:}{10011$_{2}$}{11110$_{2}$}{100101$_{2}$}{11000$_{2}$}
\mcq{479}{5 in binary syystem is:}{10$_{2}$}{101$_{2}$}{11$_{2}$}{None of these}
\mcq{480}{8 in binary system is:}{100$_{2}$}{101$_{2}$}{11$_{2}$}{None of these}
\mcq{481}{Concersion of (7)10 equivalent to binary number gives:}{110$_{2}$}{112$_{2}$}{111$_{2}$}{101$_{2}$}
\mcq{482}{19 in binary system is:}{10010$_{2}$}{11101$_{2}$}{10111$_{2}$}{10011$_{2}$}
\mcq{483}{Hexadecimal number system is based on:}{Two digits}{Ten digits}{Eight digits}{Sixteen digits}
\mcq{484}{In decimal (101)$_{2}$+(11)$_{2}$ is equal to:}{2}{8}{4}{6}
\mcq{485}{(11)$_{2}$x (11)$_{2}$ is equal to:}{1000$_{2}$}{1001$_{2}$}{1011$_{2}$}{1101$_{2}$}
\mcq{486}{$(1010)_{2}\div(10)_{2}$ is equal to:}{101$_{2}$}{10$_{2}$}{11$_{2}$}{100$_{2}$}
\mcq{487}{$(100)_{2}\div(10)_{2}$=......:}{11$_{2}$}{10$_{2}$}{101$_{2}$}{100$_{2}$}
\mcq{488}{Commonly used number system is:}{Binary number system}{Decimal number system}{Octal number system}{Hexadecimal number system}
\mcq{489}{Binary number system is based in the digits:}{0 to 9}{0 to 1}{0 to 7}{0 to 15}
\mcq{490}{Number system used in computer system is:}{Decimal}{Binary}{Octal}{Hexadecimal}

\end{document}
